\chapter{Auxiliary background}\label{sec:: ap: categorical_background}

This appendix collects a small amount of background used repeatedly in Chapters~2--5. The goal is not to give a self-contained treatment of these topics, but to gather the constructions and lemmas that are invoked later while keeping the main chapters focused on the new material of the thesis.

\section{Categories of elements and finality}

\begin{definition}\label{ap:: def: cat of els}
    Let $\cC$ be a small category and let $\cF \maps \cC \longra \set$ be a functor. The \define{category of elements} of $\cF$, denoted
    \begin{equation}
        \int \cF,
    \end{equation}
    is the category whose:
    \begin{enumerate}
        \item objects are pairs $(c,x)$ with $c\in \cC$ and $x\in \cF(c)$;
        \item morphisms $(c,x)\to (d,y)$ are morphisms $f\maps c\to d$ in $\cC$ satisfying $\cF(f)(x)=y$.
    \end{enumerate}
    There is a tautological forgetful functor
    \begin{equation}
        \int \cF \longra \cC.
    \end{equation}
\end{definition}

\begin{remark}
    If $K\maps \bDelta^{\op}\longra \set$ is a simplicial set, then its category of simplices is precisely the category of elements of $K$. This is the construction used in Chapter~2.
\end{remark}

\begin{definition}\label{ap:: def: initial functor}
    Let $K\maps \cF \longra \cD$ be a functor.
    \begin{enumerate}
        \item We say that $K$ is \define{initial} if for every complete category $\cE$ and every functor $F\maps \cD \longra \cE$, the canonical map
        \begin{equation}
            \Limit[\cD]{F} \longra \Limit[\cF]{F\circ K}
        \end{equation}
        is an isomorphism whenever either side exists.
        \item We say that $K$ is \define{final} if for every cocomplete category $\cE$ and every functor $F\maps \cD \longra \cE$, the canonical map
        \begin{equation}
            \Colimit[\cF]{F\circ K} \longra \Colimit[\cD]{F}
        \end{equation}
        is an isomorphism whenever either side exists.
    \end{enumerate}
\end{definition}

\begin{proposition}
    A functor $K\maps \cF \longra \cD$ is initial if, for every object $d\in \cD$, the comma category $d/K$ is non-empty and connected. Dually, $K$ is final if, for every object $d\in \cD$, the comma category $K/d$ is non-empty and connected.
\end{proposition}

\begin{remark}
    The homotopical version of this criterion replaces connectedness by weak contractibility. This is the guiding principle behind several contractibility arguments in Chapter~2.
\end{remark}

\section{Kan extensions}\label{ap:: subsec: Kan extensions}

\begin{definition}
    Let $F\maps \cF \longra \cE$ and $K\maps \cF \longra \cD$ be functors. A \define{right Kan extension} of $F$ along $K$ is a functor
    \begin{equation}
        \rkan_{K}(F)\maps \cD \longra \cE
    \end{equation}
    together with a natural transformation
    \begin{equation}
        \epsilon \maps \rkan_{K}(F)\circ K \Rightarrow F
    \end{equation}
    which is terminal among such pairs. Dually, a \define{left Kan extension} of $F$ along $K$ is initial among functors out of $\cD$ equipped with a natural transformation from $F$.
\end{definition}

\begin{theorem}\label{ap:: thm: right Kan extension pointwise}
    Assume that $\cE$ is complete. Then the right Kan extension of $F$ along $K$ exists and is computed pointwise by
    \begin{equation}
        \rkan_{K}(F)(d)\cong \Limit[d/K]{\big(d/K \longra \cF \xra{\,\,F\,\,} \cE\big)}.
    \end{equation}
    Dually, if $\cE$ is cocomplete, then left Kan extensions exist and are computed by colimits over the corresponding comma categories.
\end{theorem}

\section{Discrete fibrations and presheaves}

Fibrations enter the thesis in two different ways: first as the ordinary categorical shadow of presheaves, and later as the $(\oo,1)$-categorical language in which tangential and defect data are most naturally expressed.

\begin{definition}
    A functor $p\maps \cE \longra \cB$ is a \define{discrete right fibration} if for every morphism $f\maps x \to y$ in $\cB$ and every object $\widetilde{y}\in \cE$ with $p(\widetilde{y})=y$, there exists a unique morphism $\widetilde{f}\maps \widetilde{x}\to \widetilde{y}$ in $\cE$ such that $p(\widetilde{f})=f$.

    \medskip

    \nid Dually, $p$ is a \define{discrete left fibration} if $p^{\op}\maps \cE^{\op}\longra \cB^{\op}$ is a discrete right fibration.
\end{definition}

\begin{construction}\label{ap:: cons: discrete fibration}
    Let $p\maps \cE \longra \cB$ be a discrete right fibration. We define a presheaf
    \begin{equation}
        \mathsf{Un}(p)\maps \cB^{\op} \longra \set
    \end{equation}
    by sending an object $b\in \cB$ to the fibre
    \begin{equation}
        \mathsf{Un}(p)(b)=\{\widetilde{b}\in \cE : p(\widetilde{b})=b\}.
    \end{equation}
    Given a morphism $f\maps b\to b'$ in $\cB$, the map
    \begin{equation}
        \mathsf{Un}(p)(f)\maps \mathsf{Un}(p)(b') \longra \mathsf{Un}(p)(b)
    \end{equation}
    sends $\widetilde{b'}$ to the source of the unique lift of $f$ with target $\widetilde{b'}$.

    \medskip

    \nid Conversely, if $\cF\maps \cB^{\op}\longra \set$ is a presheaf, then the forgetful functor
    \begin{equation}
        \int \cF \longra \cB
    \end{equation}
    is a discrete right fibration. These two constructions are inverse to one another up to canonical isomorphism.
\end{construction}

\begin{proposition}
    The construction of \Cref{ap:: cons: discrete fibration} defines an equivalence of categories
    \begin{equation}
        \mathsf{dRFib}_{\cB} \,\, \simeq \,\, \Fun(\cB^{\op},\set).
    \end{equation}
\end{proposition}

\begin{definition}
    Let $p\maps \cE \longra \cC$ be a map of quasicategories. We say that $p$ is a \define{right fibration} if it has the right lifting property with respect to every inner horn inclusion
    \begin{equation}
        \Lambda^{n}_{i}\hookra \Delta^{n},
    \end{equation}
    for $0<i\leq n$.
\end{definition}

\begin{theorem}\label{ap:: thm: straightening unstraightening}
    For every $(\oo,1)$-category $\cC$, straightening and unstraightening induce an equivalence of $(\oo,1)$-categories
    \begin{equation}
        \mathsf{RFib}_{\cC} \,\, \simeq \,\, \fun(\cC^{\op},\mathsf{Kan}).
    \end{equation}
\end{theorem}

\begin{remark}
    Here $\mathsf{Kan}$ is a convenient model for the $(\oo,1)$-category of spaces. Restricting to right fibrations with discrete fibres recovers set-valued presheaves. This is the perspective used in Chapter~5 to encode discrete defect data.
\end{remark}

\section{A homotopy covering criterion}

\begin{theorem}\label{ap:: thm: lurie Seifert}
    Let $X$ be a topological space, let $K$ be a small category, and let
    \begin{equation}
        F\maps K \longra \mathcal{U}(X)
    \end{equation}
    be a functor to the poset of open subsets of $X$. Assume that the opens $\{F(k)\}_{k\in K}$ cover $X$ and that, for every point $x\in X$, the full subcategory
    \begin{equation}
        K_{x}=\{k\in K:\, x\in F(k)\}
    \end{equation}
    has weakly contractible nerve. Then the canonical map
    \begin{equation}
        \hoColimit[K]{F} \longra X
    \end{equation}
    is a weak homotopy equivalence.
\end{theorem}

\begin{remark}
    In Chapter~2 we use this criterion to identify certain nerves with contractible spaces by exhibiting them as homotopy colimits of contractible open sets in a contractible ambient space.
\end{remark}

\section{Differential topology}

The bordism constructions of Chapters~3 and~4 repeatedly use the same small package of differential-topological facts: proper submersions are locally trivial, bundles over contractible bases are trivial, and the same remains true in charts adapted to a boundary. We gather those statements here.

\begin{definition}\label{ap:: def: smooth fibration}
    A smooth map $\pi \maps E \longra B$ is a \define{smooth fibre bundle} if every point $b\in B$ admits an open neighbourhood $U\subseteq B$, a smooth manifold $F$, and a diffeomorphism
    \begin{equation}
        \varphi_{U}\maps \pi^{-1}(U) \xra{\,\,\cong\,\,} U\times F
    \end{equation}
    commuting with the projections to $U$.
\end{definition}

\begin{proposition}\label{ap:: prop: triviality contractibility}
    Let $\pi\maps E \longra B$ be a smooth fibre bundle with fibre $F$. If $B$ is contractible, then $\pi$ is isomorphic, as a smooth fibre bundle, to the trivial bundle
    \begin{equation}
        B\times F \longra B.
    \end{equation}
\end{proposition}
\begin{proof}
    This is standard. One way to prove it is to use path lifting to transport the fibre over a chosen basepoint along a contraction of $B$. See \cite[Corollary~10.3]{Hu94}.
\end{proof}

\begin{definition}\label{ap:: def: proper map}
    A map $f\maps M \longra N$ is \define{proper} if the inverse image of every compact subset of $N$ is compact in $M$.
\end{definition}

\begin{proposition}[Ehresmann fibration theorem]\label{ap:: prop: Ehresmann fibration no boundary}
    Let $f\maps M \longra N$ be a smooth proper submersion between smooth manifolds without boundary. Then $f$ is a smooth fibre bundle.
\end{proposition}

\begin{cor}
    Under the hypotheses of \Cref{ap:: prop: Ehresmann fibration no boundary}, all fibres of $f$ are diffeomorphic.
\end{cor}

\begin{proposition}\label{ap:: prop: ehresmann boundary}
    Let $M$ and $N$ be smooth manifolds with boundary, and let $f\maps M \longra N$ be a smooth map such that:
    \begin{enumerate}
        \item $f$ is a proper submersion;
        \item $f(\partial M)\subseteq \partial N$.
    \end{enumerate}
    Then $f$ is a smooth fibre bundle in charts adapted to the boundary. In particular, its fibres are smooth manifolds without boundary and are locally constant up to diffeomorphism.
\end{proposition}
\begin{proof}
    Away from $\partial N$ this is the classical Ehresmann theorem. Near a boundary point of $N$, choose a chart modelled on $\RR_{\geq 0}\times \RR^{n-1}$. The inward-pointing vector field in the normal direction can be lifted locally along the submersion $f$. Because $f(\partial M)\subseteq \partial N$, the lift can be chosen inward-pointing along $\partial M$, so its flow remains inside $M$. Integrating these local lifts yields local trivializations compatible with the boundary, exactly as in the standard proof of Ehresmann's theorem.
\end{proof}

\begin{remark}
    This is the version used in Chapter~3 when one studies projections to intervals in the presence of collars. Chapter~4 proves the corresponding conically smooth stratified analogue.
\end{remark}
